\subsection{\texorpdfstring{性质 \(\operatorname{Hom}_{\mathbf{A}}(\mathbf{E}, \mathbf{F})\) 关于正合序列}{性质 \textbackslash operatorname\{Hom\}\_\{\textbackslash mathbf\{A\}\}(\textbackslash mathbf\{E\}, \textbackslash mathbf\{F\}关于正合序列的相对性}}

定理1. 设A是一个环，E'、E、E''是三个A-模，u: E' → E，v: E → E''是两个同态。序列

\[
\mathrm{E} ^ {\prime} \stackrel {u} {\longrightarrow} \mathrm{E} \stackrel {v} {\longrightarrow} \mathrm{E} ^ {\prime \prime} \longrightarrow 0\tag{1}
\]

为正合的，必须且只需对于每个A-模F，序列

\[
0 \longrightarrow \operatorname{Hom} (\mathrm{E} ^ {\prime \prime}, \mathrm{F}) \stackrel {v} {\longrightarrow} \operatorname{Hom} (\mathrm{E}, \mathrm{F}) \stackrel {\bar {u}} {\longrightarrow} \operatorname{Hom} (\mathrm{E} ^ {\prime}, \mathrm{F})\tag{2}
\]

(其中 \(\bar{u} = \mathrm{Hom}(u, 1_{\mathbf{F}})\) , \(\bar{v} = \mathrm{Hom}(v, 1_{\mathbf{F}})\) ）是正合的。

假设序列(1)是正合的。如果 \(w \in \operatorname{Hom}(\mathbf{E}^{\prime \prime}, \mathbf{F})\) 且 \(\bar{v}(w) = w \circ v = 0\) ，则 \(w = 0\) ，因为 \(v\) 是满射的。因此序列（2）在 \(\operatorname{Hom}(\mathbf{E}^{\prime \prime}, \mathbf{F})\) 处是正合的。我们证明它在 \(\operatorname{Hom}(\mathbf{E}, \mathbf{F})\) 处也是正合的。由 \(\bar{u} \circ \bar{v} = \operatorname{Hom}(v \circ u, 1_{\mathbf{F}})\) （§1，第2号，公式（10））以及由序列(1)在E处正合所得的 \(v \circ u = 0\) ，得 \(\bar{u} \circ \bar{v} = 0\) ，即 \(\operatorname{Im}(\bar{v}) \subset \operatorname{Ker}(\bar{u})\) 。另一方面，如果 \(w \in \operatorname{Ker}(\bar{u})\) ，则 \(w \circ u = 0\) ，因此 \(\operatorname{Ker}(w) \supset \operatorname{Im}(u)\) 。但由于序列(1)在E处是正合的， \(\operatorname{Im}(u) = \operatorname{Ker}(v)\) ，因此 \(\operatorname{Ker}(w) \supset \operatorname{Ker}(v)\) ；由于 \(v\) 是满射的，由§1第3号注记可知存在 \(w' \in \operatorname{Hom}(\mathbf{E}^{\prime \prime}, \mathbf{F})\) ，使得

\begin{enumerate}
\def\labelenumi{(\arabic{enumi})}
\setcounter{enumi}{3}
\tightlist
\item
\end{enumerate}

\(w = w' \circ v = \bar{v}(w')\) 。因此 \(\operatorname{Ker}(\bar{u}) \subset \operatorname{Im}(\bar{v})\) ，这就完成了序列(2)正合的证明。

反之，假设（2）对每个A-模F都是正合的。由于 \(\bar{u} \circ \bar{v} = \operatorname{Hom}(v \circ u, 1_{\mathbf{F}}) = 0\) , \(w \circ v \circ u = 0\) 对每个同态 \(w: \mathrm{E}'' \to \mathrm{F}\) 成立。取 \(\mathrm{F} = \mathrm{E}''\) 和 \(w = 1_{\mathbf{E}''}\) ，首先可以看出 \(v \circ u = 0\) ，因此 \(u(\mathrm{E}') \subset \operatorname{Ker}(v)\) 。现在取 \(\mathrm{F} = \operatorname{Coker}(u)\) ，并设 \(\phi: \mathrm{E} \to \mathrm{F} = \mathrm{E}/u(\mathrm{E}')\) 为典范映射。则由定义有 \(\bar{u}(\phi) = \phi \circ u = 0\) ，因此存在 \(\psi \in \operatorname{Hom}(\mathrm{E}''\) ，F）使得 \(\phi = \bar{v}(\psi) = \psi \circ v\) ；这显然意味着 \(u(\mathrm{E}') = \operatorname{Ker}(\phi) \supset \operatorname{Ker}(v)\) ，这就证明了序列(1)在E处是正合的。最后，设 \(\theta\) 为 \(\mathrm{E}''\) 到 \(\mathrm{F} = \mathrm{E}''/v(\mathrm{E})\) 上的典范同态；于是 \(\bar{v}(\theta) = \theta \circ v = 0\) ，因此 \(\theta = 0\) ；从而 \(\mathrm{F} = \{0\}\) 且 \(v\) 是满射的。因此序列（1）在 \(\mathrm{E}''\) 处是正合的。

推论。对于A-线性映射 \(u: \mathbf{E} \to \mathbf{F}\) 为满射（相应地，双射，相应地，零映射），当且仅当对每个A-模 \(\mathbf{G}\) ，映射 \(\operatorname{Hom}(u, 1_{\mathbf{G}}): \operatorname{Hom}(\mathbf{F}, \mathbf{G}) \to \operatorname{Hom}(\mathbf{E}, \mathbf{G})\) 是单射（相应地，双射，相应地，零映射）。

只需将定理1应用于以下情形即可： \(\mathbf{E}'' = \{0\}\) （相应地 \(\mathbf{E}' = \{0\}\) ，相应地 \(\mathbf{E}'' = \mathbf{E}\) 且 \(v = 1_{\mathbb{E}}\) ).

注意，从正合序列出发

\[
0 \longrightarrow \mathrm{E} ^ {\prime} \stackrel {u} {\longrightarrow} \mathrm{E} \stackrel {v} {\longrightarrow} \mathrm{E} ^ {\prime \prime} \longrightarrow 0
\]

相应的序列

\[
0 \longrightarrow \operatorname{Hom} (\mathrm{E} ^ {\prime \prime}, \mathrm{F}) \stackrel {\bar {v}} {\longrightarrow} \operatorname{Hom} (\mathrm{E}, \mathrm{F}) \stackrel {\bar {u}} {\longrightarrow} \operatorname{Hom} (\mathrm{E} ^ {\prime}, \mathrm{F}) \longrightarrow 0
\]

不一定是正合的，换句话说，同态 \(\bar{u}\) 不一定是满射的。如果 \(E'\) 与E的一个子模等同，这就意味着从 \(E'\) 到F的线性映射并不总能扩张为从E到F的线性映射（习题11和12）。然而：

命题1. 若线性映射的正合序列

\[
0 \longrightarrow \mathrm{E} ^ {\prime} \stackrel {u} {\longrightarrow} \mathrm{E} \stackrel {v} {\longrightarrow} \mathrm{E} ^ {\prime \prime} \longrightarrow 0\tag{3}
\]

分裂（换言之，若 \(u(\mathbf{E}')\) 是 \(\mathbf{E}\) 的一个直因子），则序列

\[
0 \longrightarrow \operatorname{Hom} (\mathrm{E} ^ {\prime \prime}, \mathrm{F}) \stackrel {\bar {v}} {\longrightarrow} \operatorname{Hom} (\mathrm{E}, \mathrm{F}) \stackrel {\bar {u}} {\longrightarrow} \operatorname{Hom} (\mathrm{E} ^ {\prime}, \mathrm{F}) \longrightarrow 0
\]

是正合的且分裂。反之，若对每一个 \(A\) -模 \(F\) ，序列(4)是正合的，则序列(3)分裂。

若正合序列（3）分裂，则存在一个线性收缩 \(u': \mathbf{E} \to \mathbf{E}'\) 与 \(u\) 相关联（§1，第9号，命题15）；如果

\[
\bar {u} ^ {\prime} = \operatorname{Hom} (u ^ {\prime}, 1 _ {\mathrm{F}}): \operatorname{Hom} (\mathrm{E} ^ {\prime}, \mathrm{F}) \rightarrow \operatorname{Hom} (\mathrm{E}, \mathrm{F}),
\]

\(u' \circ u\) 是恒等映射这一事实蕴含 \(\bar{u} \circ \bar{u}'\) 是恒等映射（§1，第2号，公式（10）），因此第一个断言由§1，第9号，命题15得出。反之，假设序列（4）对 \(F = E'\) 是正合的。则存在一个元素 \(f \in \operatorname{Hom}(E, E')\) 使得 \(f \circ u = 1_{E'}\) ，结论由§1，第9号，命题15得出。

注意命题1的第一个断言也可以视为§1，第6号，命题6的推论1的一个特例，只要将 \(\operatorname{Hom}(\mathbf{E}',\mathbf{F}) \oplus \operatorname{Hom}(\mathbf{E}''\) , \(\mathbf{F}\) ）与 \(\operatorname{Hom}(\mathbf{E}' \oplus \mathbf{E}''\) , \(\mathbf{F}\) ）通过Z-线性映射 \(\operatorname{Hom}(p',\mathbf{l}_{\mathbf{F}}) + \operatorname{Hom}(p'',\mathbf{l}_{\mathbf{F}})\) 典范地等同起来，其中 \(p':\mathbf{E}' \oplus \mathbf{E}'' \to \mathbf{E}'\) 和 \(p'':\mathbf{E}' \oplus \mathbf{E}'' \to \mathbf{E}''\) 是典范投影。

定理2. 设A是一个环，F'，F，F''是三个A-模，u: F' → F，v: F → F''是两个同态。序列

\[
0 \longrightarrow \mathrm{F} ^ {\prime} \stackrel {v} {\longrightarrow} \mathrm{F} \stackrel {u} {\longrightarrow} \mathrm{F} ^ {\prime \prime}\tag{5}
\]

正合的充分必要条件是：对每个A-模E，序列

\[
0 \longrightarrow \operatorname{Hom} (\mathrm{E}, \mathrm{F} ^ {\prime}) \stackrel {\bar {u}} {\longrightarrow} \operatorname{Hom} (\mathrm{E}, \mathrm{F}) \stackrel {\bar {v}} {\longrightarrow} \operatorname{Hom} (\mathrm{E}, \mathrm{F} ^ {\prime \prime})\tag{6}
\]

(其中 \(\bar{u} = \operatorname{Hom}(1_{\mathbf{E}}, u)\) , \(\bar{v} = \operatorname{Hom}(1_{\mathbf{E}}, v)\) ）是正合的。

假设序列（5）是正合的。首先注意

\[
\bar {v} \circ \bar {u} = \operatorname{Hom} (1 _ {\mathrm{E}}, v \circ u) = 0
\]

（II，§1，第2号，公式（10）），因为 \(v \circ u = 0\) 。因此 Hom(E, F') 在 \(\bar{u}\) 下的像包含在核 N （即 \(\bar{v}\) 的核）中；令 \(f\) 是Z-模Hom(E, F')到N的同态，其图等于 \(\bar{u}\) 的图；需要证明 \(f\) 是双射，从而需要定义一个映射 \(g: \mathrm{N} \to \mathrm{Hom}(\mathrm{E}, \mathrm{F}')\) 使得 \(f \circ g\) 和 \(g \circ f\) 是恒等映射。为此，令 \(w\) 是 N 的一个元素，即一个线性映射 \(w: \mathrm{E} \to \mathrm{F}\) ，满足 \(v \circ w = 0\) 。由假设，后一关系等价于 \(w(\mathrm{E}) \subset \mathrm{Ker}(v) = u(\mathrm{F}')\) ；因此，由于 \(u\) 是单射的，存在且仅存在一个线性映射 \(w': \mathrm{E} \to \mathrm{F}'\) 使得 \(w = u \circ w'\) ，我们取 \(g(w) = w'\) ；立即可以验证 \(g\) 满足所需的条件。

反之，假设序列(6)对每个A-模E都是正合的。由于 \(\operatorname{Hom}(1_{\mathbf{E}}, v \circ u) = \bar{v} \circ \bar{u} = 0\) ，则 \(v \circ u \circ w = 0\) 对每个同态 \(w: E \to F'\) 成立。取 \(E = F'\) 和 \(w = 1_{F'}\) ，首先可以看出 \(v \circ u = 0\) ，因此 \(u(F') \subset \operatorname{Ker}(v)\) 。现在我们取 \(E = \operatorname{Ker}(v)\) ，并设 \(\phi: E \to F\) 为典范单射。则由定义有 \(\bar{v}(\phi) = v \circ \phi = 0\) ，因此存在 \(\psi \in \operatorname{Hom}(E, F')\) 使得 \(\phi = \bar{u}(\psi) = u \circ \psi\) ，这显然蕴含 \(\operatorname{Ker}(v) \subset u(F')\) ，从而完成了(5)在F处正合性的证明。最后，如果 \(\theta\) 是 \(\operatorname{Ker} u\) 的恒等映射，则 \(\bar{u}(\theta) = 0\) ，因此 \(\theta = 0\) 且 \(\operatorname{Ker} u = \{0\}\) ，这证明了(5)在 \(F'\) 处的正合性。

注记。(1) 定理2使我们能够，对于每个子模 \(\mathbf{F}'\) （属于 \(\mathbf{F}\) ），将 \(\operatorname{Hom}(\mathbf{E}, \mathbf{F}')\) 等同于 \(\mathbf{Z}\) -模 \(\operatorname{Hom}(\mathbf{E}, \mathbf{F})\) 的一个子模。当作出这一等同后，对于每个子模族 \((\mathbf{M}_{\lambda})\) （这些子模都属于 \(\mathbf{F}\) ），有

\[
\operatorname{Hom} \left(\mathrm{E}, \bigcap_ {\lambda} \mathrm{M} _ {\lambda}\right) = \bigcap_ {\lambda} \operatorname{Hom} (\mathrm{E}, \mathrm{M} _ {\lambda})
\]

，因为若 \(u \in \operatorname{Hom}(\mathbf{E}, \mathbf{F})\) 属于每一个 \(\operatorname{Hom}(\mathbf{E}, \mathbf{M}_{\lambda})\) ，那么对所有 \(x \in \mathbf{E}\) , \(u(x) \in \mathbf{M}_{\lambda}\) 对所有 \(\lambda\) 成立，因此 \(u\) 把 \(\mathbf{E}\) 映到 \(\bigcap_{\lambda} \mathbf{M}_{\lambda}\) 内.

推论。对于A-线性映射 \(u: \mathbf{E} \to \mathbf{F}\) ，为使它是单射的，必须且只需对每个A-模 \(\mathbf{G}\) ，映射 \(\operatorname{Hom}(1_{\mathbf{G}}, u): \operatorname{Hom}(\mathbf{G}, \mathbf{E}) \to \operatorname{Hom}(\mathbf{G}, \mathbf{F})\) 是单射的。

只需将定理2应用于以下情形： \(\mathbf{F}' = \{0\}\) .

从一个正合序列出发

\[
0 \longrightarrow \mathrm{F} ^ {\prime} \stackrel {u} {\longrightarrow} \mathrm{F} \stackrel {v} {\longrightarrow} \mathrm{F} ^ {\prime \prime} \longrightarrow 0
\]

相应的序列

\[
0 \longrightarrow \operatorname{Hom} (\mathrm{E}, \mathrm{F} ^ {\prime}) \stackrel {\bar {u}} {\longrightarrow} \operatorname{Hom} (\mathrm{E}, \mathrm{F}) \stackrel {\bar {v}} {\longrightarrow} \operatorname{Hom} (\mathrm{E}, \mathrm{F} ^ {\prime \prime}) \longrightarrow 0
\]

不一定是正合的，换句话说， \(\bar{v}\) 不一定是满射的。如果 \(F'\) 被等同于F的一个子模，且 \(F''\) 被等同于商模 \(F/F'\) ，这就意味着从E到 \(F''\) 的线性映射不一定具有形式 \(v \circ w\) ，其中 w 是从E到F的线性映射。然而：

命题 2. 若正合序列

\[
0 \longrightarrow \mathrm{F} ^ {\prime} \stackrel {u} {\longrightarrow} \mathrm{F} \stackrel {v} {\longrightarrow} \mathrm{F} ^ {\prime \prime} \longrightarrow 0\tag{7}
\]

分裂（换言之，若 \(u(\mathbf{F}')\) 是 \(\mathbf{F}\) 的一个直因子），则序列

\[
0 \longrightarrow \operatorname{Hom} (\mathrm{E}, \mathrm{F} ^ {\prime}) \stackrel {\bar {u}} {\longrightarrow} \operatorname{Hom} (\mathrm{E}, \mathrm{F}) \stackrel {\bar {v}} {\longrightarrow} \operatorname{Hom} (\mathrm{E}, \mathrm{F} ^ {\prime \prime}) \longrightarrow 0\tag{8}
\]

是正合的且分裂。反之，若序列(8)对每个A-模E都是正合的，则正合序列(7)分裂。

第一个断言由以下事实得出：

\[
\operatorname{Hom} (\mathbf {E}, \mathbf {F} ^ {\prime}) \oplus \operatorname{Hom} (\mathbf {E}, \mathbf {F} ^ {\prime \prime})
\]

被典范地等同于 \(\operatorname{Hom}(\mathbf{E}, \mathbf{F}' \oplus \mathbf{F}'')\) ，其中使用了 \(\mathbf{Z}\) -线性映射 \(\operatorname{Hom}(1_{\mathbb{E}}, j') + \operatorname{Hom}(1_{\mathbb{E}}, j'')\) , \(j': \mathbf{F}' \to \mathbf{F}' \oplus \mathbf{F}''\) 和 \(j'': \mathbf{F}'' \to \mathbf{F}' \oplus \mathbf{F}''\) 为典范嵌入（§1，第6号，命题6的推论1）。反之，若序列(8)对 \(\mathbf{E} = \mathbf{F}''\) 是正合的，则存在一个元素 \(g \in \operatorname{Hom}(\mathbf{F}''\) , \(\mathbf{F}\) ）使得 \(v \circ g = 1_{\mathbf{F}''}\) ，且结论由§1，第9号，命题15得出。

注记(2)。本号的结果对于所有带算子的交换群无需修改即可成立。

